\chapter{共识协议骨架}\label{app:protocol}

本附录勾勒第 \ref{sec:mechanism} 节论证其必要、第 \ref{sec:consensus} 节为之建模的机制：
一组独立参与方把意图转换为一份规格、他人随后可以据此执行的程序。之所以把它写下来，
是因为本章的核心主张——意图层的保证来自对抗性共识，而不是来自验证者——若不对这样一种
程序至少给出一个骨架，就是空洞的；也因为"该交叉点\emph{无人占据}"（第
\ref{ch:literature} 章）这一主张，对照一份具体方案来评估比对照一段描述更容易。

本附录有三件事不是。它不是实现；其中没有任何部分被构建过或运行过。它未经确认；本文不
报告任何实验，其中的参数被陈述为有待确定的量，而不是被测量出的值。它也不是关于正确性的
主张：下面的协议可以被严格照做，却仍然产出一份误述意图的规格，原因由命题
\ref{prop:three-grounds} 列出。它确实主张的是一个形状：本文其余部分加在意图层上的义务，
可以由一个具有这些步骤的程序来履行；而每一步在已部署实践中都有一个对应物。

\section{被审定的对象}

该协议产出一份\term{审定记录}，它是链条其余部分所消费的意图层制品。它的字段是制品证据
清单（表 \ref{tab:evidence-package}）在意图层的对应物，而这种对应是有意的：两份清单由
同一套下游机制读取。

\begin{center}
\small
\begin{longtable}{L{0.9cm}L{5.2cm}L{8.5cm}}
\caption{审定记录}\label{tab:attestation}\\
\toprule
\textbf{\#} & \textbf{字段} & \textbf{内容以及为何需要它} \\
\midrule
\endfirsthead
\toprule
\textbf{\#} & \textbf{字段} & \textbf{内容以及为何需要它} \\
\midrule
\endhead
\bottomrule
\endlastfoot
1 & 规格文本 & 形式对象 $\Phi$ 本身，连同其范围 $\Sigma$ 与其所声明的环境假设 $A$。 \\
2 & 意图陈述 & 由具备资格的一方所陈述的需求，使用该方所用的任何记号。逐字保留：转述正是
  意图丢失之处，而字段 1 与字段 2 之间的鸿沟，恰恰是命题 \ref{prop:three-grounds}(i) 说
  无法被机械闭合的那个对象。 \\
3 & 审定者身份与利益相关者类别 & 谁作出了审定，以及每一方代表何种利益。之所以需要，
  是因为总体的可靠性取决于错误来源的\emph{独立性}，而利益相关者类别是这种独立性的
  可观察代理（第 \ref{sec:mechanism} 节）。 \\
4 & 证据来源 & 对每位审定者而言，其判断所依据的来源。正是这个字段使 $q$ 在原理上可测：
  两位审定者若依赖同一个来源，则无论他们是否知道，都是相关的；这就是
  \cite{dietrich2004} 的共同原因机制。 \\
5 & 异议记录 & 被提出的反对意见，以及它们是如何被解决的。一份没有异议的记录是关于程序
  的证据，而不是关于规格的证据。 \\
6 & 阈值与规则 & 所用的判定规则、阈值及其理由。 \\
7 & 授权机构 & 谁采用了该规格，以何种资格（第 \ref{sec:authority} 节）。它与审定者
  不同：审定是判断，授权是决定。 \\
8 & 有效视界与重新审定触发条件 & 审定何时失效，以及哪些事件强制重新审定、而不是仅仅
  重新验证（式 \eqref{eq:spec-decay}）。 \\
9 & \textbf{已知未解决的意图} & 程序无法解决的需求。它是表 \ref{tab:evidence-package}
  第 9 项中未覆盖行为的对应物，也是最可能被省略的字段，因为记录它会削弱结果的外观。 \\
10 & \textbf{重新审定规则} & 规格、意图、环境或授权机构的哪些变化会使该记录失效。 \\
\end{longtable}
\end{center}

字段 9 与字段 10 在这里的分量，与第 9 项和第 10 项在那里的分量相同。一份省略字段 9 的
记录不是更强的记录；它是一份边界未知的记录，而下游证据链会把该规格当作完整的来处理，
尽管它并不完整。

\section{程序}

以下步骤按执行的先后顺序陈述。步骤 1--3 是构造性的，步骤 4 是对抗性的，步骤 5 是决定，
步骤 6--7 是维护。

\begin{enumerate}[label=\textbf{P\arabic*.},leftmargin=*]
\item \term{框定意图}。具备资格的一方陈述目标、约束与验收标准，并指明范围 $\Sigma$。
      输出：字段 2。
\item \term{推导规格}。一个或多个提出者由意图产出候选规格。这一步可以被自动化，而在
      第 \ref{ch:program} 章的框架中，这正是预期生成最便宜之处。输出：字段 1 的候选值。
\item \term{使推导可追溯}。候选规格的每一条款都链接到它在意图陈述中所履行的部分。没有
      链接的条款是未被请求的；没有对应条款的意图部分是未被表达的。这一步是使字段 9 变得
      可写、而不是停留为愿望的机制，也是近来的可追溯精化工作所自动化的那项纪律（第
      \ref{sec:gaps} 节）。
\item \term{对抗性地审定}。从有意选取为彼此不同的利益相关者类别中抽取的审定者，对照
      意图评判候选规格，并带有明确的指示：去寻找失败，而不是去认可。输出：字段 3、4
      与 5。有两项设计选择比其他选择更要紧。审定者的遴选必须依据\emph{利益与证据的
      多样性}，因为可靠性论证所要求的正是这一点，而不只是依据能力；异议必须被记录，
      而不是被消解掉，因为一份未解决的分歧记录是这一步最有信息量的输出。
\item \term{判定并授权}。对审定结果施加一条阈值规则，并由授权机构采用该规格。这两个
      行为是不同的，必须被分别记录（字段 6 与 7），理由由第 \ref{sec:authority} 节给出：
      一份规格可以是正确的而未被授权，也可以是被授权的而错误的。
\item \term{设定视界与触发条件}。该记录被赋予一个有效视界，以及一份强制重新审定的事件
      清单。视界不是装饰性的；式 \eqref{eq:spec-decay} 之所以把保证条件带上时间下标，
      正是因为一份已审定的规格可以在每一项机械检查都继续通过的同时，偏离一个移动中的
      意图。
\item \term{按触发条件重新审定}。一旦字段 10 中的任一触发条件发生，程序就对受影响的
      条款从步骤 1 重新开始。对实现的重新验证不是替代品：如果意图移动了，那个曾被验证
      的制品就不再是正确的待验证制品。
\end{enumerate}

\section{每一步买到了什么、没买到什么}

把协议写出来的价值在于，它把问题中由不同工具实际处理的部分分开了；本章的论证需要
这一点，而一段散文式描述会把它掩盖。

\emph{步骤 2 与 3 大体上可机械化}，也是自动化最能发挥作用之处。它们处理的是表达性问题。
本文没有任何内容否认它们可以被很好地自动化；框架恰恰预期如此。

\emph{步骤 4 不可机械化，也无法使之可机械化}，因为它的输入包含关于意图本应是什么的
判断。这就是需求发现问题，而协议对它的贡献是程序性的：它强制把分歧记录下来、把证据
来源声明出来，从而把一个不可测的量转换为一个可估计的量。这是所能主张的最大限度，而它
比第 \ref{sec:consensus} 节的可靠性论证若要成为一项预测所需要的还要少。

\emph{步骤 5 与步骤 1--4 是不同种类的行为}，而这就是第 \ref{sec:authority} 节加入到框架
中的要点。无论适当多样的审定者作出多少正确的判断，都不能使结果具有约束力；资格不是靠
准确性挣来的。一个只记录审定、却不记录授权的协议，对于"这份规格从来就不是任何人可以
采用的"这一反对意见，将无从回应。

\emph{步骤 6 与 7 是已部署实践中最常缺席的部分}，也是生命周期视角使之无法回避的部分。
它们的缺席在审定的那一刻是不可见的，只有在此之后才显现出来：此时意图已经移动，而记录
仍然说着它原来所说的话。

\section{协议自身的失效模式}

一份只列出步骤的协议骨架会是一份推销文件。下面这些模式是本设计最易暴露于其中的模式，
它们被陈述为程序的属性，而不是其参与者的属性。

\begin{enumerate}[label=(\alph*)]
\item \term{选择俘获}。如果审定者的集合由规格正受评判的那一方挑选，多样性就可以在名义
      上被满足，而错误仍然相关。字段 3 检测不到这一点，因为被俘获的名单仍然是名单。
      缓解手段是程序性的，且在协议之外：遴选规则必须在候选规格为人所知之前就固定下来。
\item \term{来源坍缩}。来自不同利益相关者类别的审定者，若取材于同一个来源，就会变得
      相关；而字段 4 正是使这一点可见的东西。在来源未被声明之处，协议退化为 Condorcet
      情形，其前提一般无法被同时满足 \cite{dietrich2008}。
\item \term{阈值表演}。在一个规模小且相关的名单上施加高阈值，所得到的置信度是式
      \eqref{eq:chernoff} 的界所不支持的，因为该界的有效样本量是由 $\bar\rho$ 设定的，
      而不是由计数设定的。这就是注 \ref{rmk:design-effect} 的设计效应要点以治理失效的
      形式出现。
\item \term{不可证伪的验收}。如果在步骤 1 没有指定任何可观察的后果，审定就无法在事后被
      评分，而协议产出的记录与一份正确的记录无法区分。这就是注
      \ref{rmk:consensus-status} 为整个共识层记录下的局限，也是在缺乏可评分结果的领域
      中反对依赖该协议的最强论据。
\item \term{条款静默丢失}。步骤 3 只有在意图陈述本身完整时，才能使未被表达的意图变得
      可见。在具备资格的一方要等到看见自己不想要什么才能说出自己想要什么之处，字段 9
      是从一个移动的目标写下的，而该记录的完备性主张比它的格式所暗示的更弱。
\end{enumerate}

这些模式没有一个能靠更好的阈值修复。其中三个是名单与来源的属性，而不是判定规则的属性；
第四个是领域的属性。这一分布是对共识协议能贡献什么、不能贡献什么的诚实总结：它可以使
相关性质\emph{被记录下来、因而可审计}，但它不能使它们\emph{为真}。
